\chapter{Control and State Estimation}\label{ch:control-and-state-estimation}

\section{Current Implementation Status}\label{current-implementation-status}

The inspected firmware contains types and state paths needed by a
balance controller, including pitch, pitch rate, roll, and roll rate,
but no completed balance-control module was confirmed. The current
system should therefore be described as control-ready infrastructure
rather than a demonstrated self-balancing implementation. The
architecture in this chapter is the proposed completion path and the
basis of the simulation draft. Final controller claims require the exact
source revision, gains, sample interval, actuator mode, and hardware
logs.

\par\medskip\noindent\singlespacing
\phantomsection\addcontentsline{lot}{table}{\protect\numberline{7.1}Control and estimation functions and current status}
\textbf{Table 7.1. Control and estimation functions and current status}\par

\begin{longtable}[]{@{}
  >{\raggedright\arraybackslash}p{(\columnwidth - 6\tabcolsep) * \real{0.2090}}
  >{\raggedright\arraybackslash}p{(\columnwidth - 6\tabcolsep) * \real{0.2761}}
  >{\raggedright\arraybackslash}p{(\columnwidth - 6\tabcolsep) * \real{0.2164}}
  >{\raggedright\arraybackslash}p{(\columnwidth - 6\tabcolsep) * \real{0.2985}}@{}}
\toprule\noalign{}
\begin{minipage}[b]{\linewidth}\raggedright
\textbf{Function}
\end{minipage} & \begin{minipage}[b]{\linewidth}\raggedright
\textbf{Current evidence}
\end{minipage} & \begin{minipage}[b]{\linewidth}\raggedright
\textbf{Status used in thesis}
\end{minipage} & \begin{minipage}[b]{\linewidth}\raggedright
\textbf{Completion evidence}
\end{minipage} \\
\midrule\noalign{}
\endhead
\bottomrule\noalign{}
\endlastfoot
Raw IMU acquisition & Portable ISM330DHCX driver and ESP32 SPI path &
Implemented\slash{}bring-up & Rate, missed-sample, and noise measurements \\
Mount transform and gyro calibration & Reported in firmware data path &
Implemented, final calibration unverified & Axis test and retained
calibration values \\
Complementary attitude estimate & Reported estimator path & Implemented
infrastructure & Static\slash{}dynamic comparison with reference angle \\
Balance-state snapshot & Balance types and application state reported &
Implemented interface & Atomicity and freshness test \\
Balance control law & No confirmed controller source & Planned & Code,
gains, timing, simulation and hardware results \\
Wheel wrench allocation & No confirmed mapping\slash{}roller signs & Planned &
Geometry matrix, sign test, saturation logic \\
Yaw\slash{}lateral outer loops & No confirmed closed-loop implementation &
Planned & Code and motion-tracking results \\
Safety-gated motor output & Supervisor\slash{}safety command gating reported &
Partially implemented & Fault-injection and stop-time results \\
\end{longtable}

\section{Proposed Cascaded Architecture}\label{proposed-cascaded-architecture}

The proposed controller separates stabilization, planar motion,
allocation, and actuator enforcement. The innermost balance loop uses
pitch and pitch rate, with translational velocity or position included
directly in a state-feedback law or applied through a slower outer loop.
Lateral and yaw controllers produce desired body forces or moments. An
allocation stage maps the combined body command to four wheel commands,
enforces per-actuator limits, and reserves balance authority. A
supervisor controls transitions among disarmed, ready, balancing, and
faulted states.

\begin{longtable}[]{@{}
  >{\raggedright\arraybackslash}p{(\columnwidth - 0\tabcolsep) * \real{1.0000}}@{}}
\toprule\noalign{}
\begin{minipage}[b]{\linewidth}\raggedright
\textless INSERT FIGURE 7.1: Draw reference conditioning, x\slash{}y\slash{}yaw outer
loops, pitch estimator, inner balance controller, body-wrench command,
four-wheel allocator, saturation and slew limits, safety permission, CAN
actuator interface, and feedback paths. Mark implemented versus planned
blocks\textgreater{}\\
What the figure should show: Draw reference conditioning, x\slash{}y\slash{}yaw outer
loops, pitch estimator, inner balance controller, body-wrench command,
four-wheel allocator, saturation and slew limits, safety permission, CAN
actuator interface, and feedback paths. Mark implemented versus planned
blocks\\
Likely source: firmware architecture and MATLAB controller model\strut
\end{minipage} \\
\midrule\noalign{}
\endhead
\bottomrule\noalign{}
\endlastfoot
\end{longtable}

\begin{center}\singlespacing
\phantomsection\addcontentsline{lof}{figure}{\protect\numberline{7.1}Proposed cascaded control and allocation architecture.}
\textbf{Figure 7.1. Proposed cascaded control and allocation architecture.}
\end{center}

The controller should not command position while disarmed or fallen. On
entry to balancing, integrators and reference states should be
initialized from the measured state to prevent a command step. Operator
inputs should be deadbanded, scaled, rate limited, and set to zero when
CRSF data become stale. Display and telemetry functions may observe
control state but may not own motor commands.

\section{State Estimation}\label{state-estimation}

\subsection{Calibration and frame transformation}\label{calibration-and-frame-transformation}

The IMU driver first converts raw counts to physical units using the
configured full-scale ranges. A fixed mounting matrix transforms sensor
axes into the body frame. Gyro bias is estimated only while the robot is
confirmed stationary, and the calibration record includes sample count,
duration, temperature, mean, and standard deviation. Accelerometer
offsets and scale can be checked with multiple static orientations. The
estimator marks its output invalid during initialization, sensor errors,
or stale acquisition.

\subsection{Complementary pitch estimate}\label{complementary-pitch-estimate}

Let \ensuremath{\omega}\_p be the body-frame gyro component about the pitch axis, \ensuremath{\theta}\_a the
gravity-derived pitch angle, and \ensuremath{\Delta}t the measured sample interval. A
discrete complementary filter is given by Equation (7.1).

\begin{equation}
\hat{\theta}_{k}=\alpha\left(\hat{\theta}_{k-1}+\Delta t\,\omega_{p,k}\right)
+(1-\alpha)\theta_{a,k}
\label{eq:7-1}
\end{equation}

The coefficient \ensuremath{\alpha} determines the crossover between the integrated gyro
and accelerometer. It should be computed from a documented time constant
and actual \ensuremath{\Delta}t, not copied as a unitless tuning constant without context.
During forward or lateral acceleration, the accelerometer does not
measure gravity alone; the accelerometer correction can therefore tilt
the estimate in the direction of specific force. Validation should
compare the estimated pitch with an external reference during static
tilts, manual rotations, motor acceleration, and disturbance recovery.

\begin{longtable}[]{@{}
  >{\raggedright\arraybackslash}p{(\columnwidth - 0\tabcolsep) * \real{1.0000}}@{}}
\toprule\noalign{}
\begin{minipage}[b]{\linewidth}\raggedright
\textless INSERT FIGURE 7.2: Show raw accelerometer and gyroscope
samples, unit conversion, mounting transform, bias calibration, gravity
angle, gyro integration, complementary fusion, validity\slash{}freshness
checks, and the published pitch\slash{}pitch-rate snapshot\textgreater{}\\
What the figure should show: Show raw accelerometer and gyroscope
samples, unit conversion, mounting transform, bias calibration, gravity
angle, gyro integration, complementary fusion, validity\slash{}freshness
checks, and the published pitch\slash{}pitch-rate snapshot\\
Likely source: firmware data-flow diagram and estimator
implementation\strut
\end{minipage} \\
\midrule\noalign{}
\endhead
\bottomrule\noalign{}
\endlastfoot
\end{longtable}

\begin{center}\singlespacing
\phantomsection\addcontentsline{lof}{figure}{\protect\numberline{7.2}Attitude-estimation signal path from raw IMU data to balance state.}
\textbf{Figure 7.2. Attitude-estimation signal path from raw IMU data to balance state.}
\end{center}

\section{Balance Control}\label{balance-control}

\subsection{State-feedback option}\label{state-feedback-option}

For the discrete linear model, the initial proposed law is Equation
(7.2), where x-hat contains estimated longitudinal and pitch states, r
contains motion references, K is the state-feedback gain, and N\_r
supplies reference feedforward or integral augmentation.

\begin{equation}
\mathbf{u}_{k}=-\mathbf{K}\hat{\mathbf{x}}_{k}+\mathbf{N}_{r}\mathbf{r}_{k},\qquad
\mathbf{K}=\operatorname{dlqr}(\mathbf{A}_{d},\mathbf{B}_{d},\mathbf{Q},\mathbf{R})
\label{eq:7-2}
\end{equation}

The continuous model is discretized at the measured control interval. Q
and R are selected after scaling the states to meaningful limits so that
numerical magnitude does not substitute for engineering priority. The
eigenvalues of the closed-loop discrete system, control effort,
saturation time, and robustness to parameter changes are reported.
Integral action may be added for steady velocity or position error only
after the inner balance loop is stable and anti-windup is implemented.

\subsection{PID option and anti-windup}\label{pid-option-and-anti-windup}

A cascaded PID implementation remains a valid alternative for initial
hardware tuning. The pitch controller uses angle error, measured or
filtered pitch rate, and carefully limited integral action. An outer
velocity loop modifies the pitch reference. Equation (7.3) shows a
continuous representation with back-calculation anti-windup.

\begin{equation}
u(t)=K_p e(t)+K_i\!\int e(t)\,dt+K_d\dot{e}(t),\qquad
\dot{z}=e+K_{aw}\left(u_{sat}-u\right)
\label{eq:7-3}
\end{equation}

Derivative action should use measured angular rate or a filtered
derivative rather than differentiating a noisy angle. The integral state
is reset or frozen when disarmed, fallen, or saturated in a direction
that would increase windup. The final thesis should report which
controller was implemented and use the other only as a comparison if it
was actually tested.

\section{Lateral, Yaw, and Motion Control}\label{lateral-yaw-and-motion-control}

Longitudinal force is coupled most directly to pitch. Lateral and yaw
commands are generated in body coordinates and passed to the allocator
with lower priority than the balance channel. The lateral controller may
begin as open-loop velocity command scaling because roller slip can make
encoder-only lateral speed estimation unreliable. Closed-loop lateral
velocity requires an independent velocity estimate or a validated
wheel-slip model. Yaw can be estimated from the gyro and corrected with
an external heading reference if required by the experiment.

The allocator uses Equation (6.6) with measured H. Each wheel command is
converted to the actuator\textquotesingle s selected operating mode. If
the AK40-10 is operated in servo velocity mode, the balance design must
account for the actuator\textquotesingle s internal velocity loop and
command period. If torque or MIT mode is selected, torque scaling,
stiffness, damping, and current limits must be documented. Mode
selection is not finalized in the available evidence.

\section{Saturation, Rate Limits, and Safety}\label{saturation-rate-limits-and-safety}

Command limits are applied in physical units after allocation. The
implemented order should be documented because independent clipping can
distort the desired wrench. Equation (7.4) represents a rate limit
followed by magnitude saturation; the actual design may use a
constrained allocator so that all limits are considered together.

\begin{equation}
\boldsymbol{\tau}_{cmd}=\operatorname{sat}_{\boldsymbol{\tau}_{max}}
\left(\operatorname{rateLimit}\left(\boldsymbol{\tau}^{\star},\dot{\boldsymbol{\tau}}_{max}\right)\right)
\label{eq:7-4}
\end{equation}

Safety checks include IMU validity, estimate age, pitch and pitch-rate
bounds, actuator-feedback age, CAN state, battery voltage, motor faults,
operator link state, emergency-stop state, and explicit arming intent. A
fault removes torque permission and requests a defined safe motor state.
Automatic re-arming after a fault is prohibited unless the final hazard
analysis supports it.

\begin{longtable}[]{@{}
  >{\raggedright\arraybackslash}p{(\columnwidth - 0\tabcolsep) * \real{1.0000}}@{}}
\toprule\noalign{}
\begin{minipage}[b]{\linewidth}\raggedright
\textless INSERT FIGURE 7.3: Show DISARMED, INITIALIZING, READY,
BALANCING, FALLEN, and FAULT states with guarded transitions for
calibration completion, arm request, valid sensors, pitch window, stale
data, motor fault, emergency stop, and manual reset\textgreater{}\\
What the figure should show: Show DISARMED, INITIALIZING, READY,
BALANCING, FALLEN, and FAULT states with guarded transitions for
calibration completion, arm request, valid sensors, pitch window, stale
data, motor fault, emergency stop, and manual reset\\
Likely source: supervisor state-machine implementation\strut
\end{minipage} \\
\midrule\noalign{}
\endhead
\bottomrule\noalign{}
\endlastfoot
\end{longtable}

\begin{center}\singlespacing
\phantomsection\addcontentsline{lof}{figure}{\protect\numberline{7.3}Safety-oriented operating-state transitions.}
\textbf{Figure 7.3. Safety-oriented operating-state transitions.}
\end{center}

\par\medskip\noindent\singlespacing
\phantomsection\addcontentsline{lot}{table}{\protect\numberline{7.2}Controller and estimator parameters requiring final values}
\textbf{Table 7.2. Controller and estimator parameters requiring final values}\par

\begin{longtable}[]{@{}
  >{\raggedright\arraybackslash}p{(\columnwidth - 6\tabcolsep) * \real{0.2941}}
  >{\raggedright\arraybackslash}p{(\columnwidth - 6\tabcolsep) * \real{0.1471}}
  >{\raggedright\arraybackslash}p{(\columnwidth - 6\tabcolsep) * \real{0.1765}}
  >{\raggedright\arraybackslash}p{(\columnwidth - 6\tabcolsep) * \real{0.3824}}@{}}
\toprule\noalign{}
\begin{minipage}[b]{\linewidth}\raggedright
\textbf{Parameter}
\end{minipage} & \begin{minipage}[b]{\linewidth}\raggedright
\textbf{Unit}
\end{minipage} & \begin{minipage}[b]{\linewidth}\raggedright
\textbf{Final value}
\end{minipage} & \begin{minipage}[b]{\linewidth}\raggedright
\textbf{Selection\slash{}verification method}
\end{minipage} \\
\midrule\noalign{}
\endhead
\bottomrule\noalign{}
\endlastfoot
Estimator sample period & s & \textless INSERT\textgreater{} & Measured
timestamps and configured IMU ODR \\
Control sample period & s & \textless INSERT\textgreater{} &
Schedulability and timing measurement \\
Complementary time constant / \ensuremath{\alpha} & s / --- &
\textless INSERT\textgreater{} & Noise and external-angle comparison \\
LQR Q and R or PID gains & various & \textless INSERT\textgreater{} &
Document scaling, simulation, restrained tuning, final code revision \\
Torque\slash{}current limit per actuator & N\ensuremath{\cdot}m / A &
\textless INSERT\textgreater{} & Installed motor configuration and
thermal\slash{}current tests \\
Command slew limit & N\ensuremath{\cdot}m\slash{}s or mode-specific &
\textless INSERT\textgreater{} & Avoid current spikes while preserving
recovery \\
Arm pitch window & rad & \textless INSERT\textgreater{} & Safe test
procedure and recovery capability \\
Fall pitch threshold & rad & \textless INSERT\textgreater{} & Mechanical
clearance and loss-of-recovery evidence \\
Sensor\slash{}feedback timeout & s & \textless INSERT\textgreater{} &
Worst-case timing and fault-injection test \\
\end{longtable}

\section{Stability and Validation}\label{stability-and-validation}

For a linear controller, local stability is assessed from the
closed-loop poles or eigenvalues of the discrete system. This analysis
is necessary but insufficient because saturation, rate limiting, delay,
estimator error, and slip are nonlinear. Simulation should sweep mass,
center-of-mass height, inertia, actuator delay, friction, and
battery-dependent limits. Hardware validation then begins with low
torque and a restraint, advances to brief free-standing trials, and
records every fall or fault.

\begin{longtable}[]{@{}
  >{\raggedright\arraybackslash}p{(\columnwidth - 0\tabcolsep) * \real{1.0000}}@{}}
\toprule\noalign{}
\begin{minipage}[b]{\linewidth}\raggedright
\textless INSERT RESULT: Final controller type, equations as
implemented, exact state vector, sample interval, gains, filters,
reference limits, saturation logic, stability analysis, simulation
response, and hardware response. Cite the Git commit and configuration
file containing each parameter.\textgreater{}
\end{minipage} \\
\midrule\noalign{}
\endhead
\bottomrule\noalign{}
\endlastfoot
\end{longtable}
